\section{クロックとCPIの重み付き計算}
\label{theory:information:cpu}
命令種類\(j\)の実行数\(N_j\)、平均サイクル数\(\mathrm{CPI}_j\)、周波数\(f\)ならCPU実行時間は\(\sum_jN_j\mathrm{CPI}_j/f\)。CPIを単純平均せず命令数で重み付けする。GHzは毎秒\(10^9\)サイクルなので、サイクル数との商は秒になる。周波数だけの大小から異なる命令構成・CPUの所要時間を決めない。

\section{参照するたびにLRUの順序を更新する}
\label{theory:information:lru}
LRUは最後の参照から最も長く経過したページを置換する。主記憶を古い参照から新しい参照の順に並べ、ヒットしたページも末尾へ移す。ミスで容量が満杯なら先頭を追い出して挿入する。FIFOの挿入時刻とLRUの参照時刻は違う。
原典がミスに外部アクセス時間M、ヒットにHを割り当てるモデルなら、両回数を数えて足す。実機のページフォールトはOS処理・転送等も含み、原典の単純時間を全実装の保証としない。

\section{検査ビットと最小距離の意味}
\label{theory:information:hamming}
ハミング距離は等長二列の異なる位置数で、XORの1の数に等しい。検査ビットは情報ビットのmod2線形結合で作り、受信語が満たすべき式を与える。最小距離\(d_{\min}\)の符号は\(t\)ビット訂正に\(2t<d_{\min}\)が必要で、距離3なら高々1ビットの誤りを一意に訂正できる。
受信語自体と全一ビット反転候補を検査し、式に合う候補を選ぶ。誤り上限を超える場合には別の符号語へ誤訂正し得るので、条件を保持する。情報ビットの誤りも検査ビットの誤りも候補に含める。

\section{状態遷移の矢印とガードを読む}
\label{theory:information:state}
状態機械は、現在状態・イベント・条件から次状態を指定する。ガードは遷移を許す条件である。エレベータの上昇は現在階が目的階より小さい間継続し、等しくなれば停止する。下降は大小を逆にする。
図の条件を文字列だけで判断せず、矢印の始点と終点を確認する。同じ条件でも異なる現在状態からの遷移を表す。扉や故障等を省いた教材モデルの範囲を、実機の安全制御の全仕様へ広げない。
